\vfill%
\pagebreak
\section{Blocks}
\label{sec:flow:blocks}

A block is a sequence of instructions between braces. It opens a scope, and it
is an expression whose value is that of its last expression.

\subsection{Grammar}

\begin{lstlisting}[style=grammarVerb]
block       := '{' declaration* instruction* expression? '}'
instruction := expression ';'
             | block_expression
             | 'let' variable_declaration ';'
\end{lstlisting}

\token{declaration} is a local declaration -- a function, a type, a module --
visible from the whole block. Local declarations precede the first instruction;
a declaration after an instruction is a syntax error.\token{variable\_declaration} is specified in
Section~\ref{sec:memory:local_variable_declaration}. A
\token{block\_expression} is an expression that ends with a block: a block
itself, \token{if}, \token{loop}, \token{while}, \token{for},
\token{match}. It can be followed by another instruction without a semicolon.
Any other expression, and every \token{let}, must be followed by one.

A block with scope guards, which the grammar above omits, is specified in
Chapter~\ref{chap:error}.

\subsection{Evaluation}

The instructions of a block are evaluated in order, from the opening brace to
the closing one. The block opens a scope: a variable declared inside it lives
until the closing brace, and cannot be named after it
(Section~\ref{sec:memory:local_variable_declaration}).

\subsection{Value and type}

When the block ends with an \token{expression} not followed by a semicolon, the
value of the block is the value of that expression, and its type is the type of
that expression. Otherwise, the block has no value, and its type is
\token{void}. A semicolon discards the value of the expression it follows, so
\token{\{ 3; \}} is of type \token{void}.

\begin{lstlisting}[style=coloredverbatim, caption=Value of a block, label=lst:flow:block_value]
fn main() {
  let x = {
    let a = 3;
    a * 2
  }; // x = 6, and a no longer exists
  println(x);
}
\end{lstlisting}

The value of an expression used as an instruction is discarded, whatever its
type. \token{n + 1;} and \token{if c \{ 1 \} else \{ 2 \}} are valid
instructions.

A block whose evaluation always reaches a breaking instruction
(Section~\ref{sec:flow:unreachable}) has no value, and that is not an error when
the block is a branch of a construct that has one.

\subsection{Diagnostics}

\begin{center}
  \begin{adjustbox}{max width=\linewidth}
    \begin{tabular}{|l|l|}
      \hline
      Code & Raised when\\[0pt]
      \hline
      \hline
      E4039 & a variable is declared with a value of type \token{void}, such as a block with no value\\[0pt]
      E4260 & an instruction follows a breaking instruction in the same block\\[0pt]
      \hline
    \end{tabular}
  \end{adjustbox}
\end{center}
